\subsection{Deriva por familia y control hasta la primera salida}\label{macro:control}

\begin{lema}[Pérdida condicional de un patrón]\label{macro:deriva}
Supongamos que el padre cumple los umbrales de la Proposición~\ref{macro:maximo} y que $k\ge\delta n$.
Para un conjunto $\mathcal L$ de un número acotado de líneas distintas disponibles, sea $H^*(\mathcal
L)$ la suma de un uno por fila o columna y de $\ell_f^*$ por cada diagonal o antidiagonal $f$. En un
paso,
\[
 P^{\rm B}(\text{se ocupa alguna línea de }\mathcal L\mid\mathcal F_s)
 =\frac{H^*(\mathcal L)}{k}
       +O_\delta\!\left(\frac{\zeta+k^{-1}}{k}\right).
\]
Las tasas de las fibras son las de la Definición~\ref{macro:fibras}, y la del término rectangular $T$
es
\[
 H_T^*=4+\sum_{m\in T}\bigl(\ell_{D(e_m)}^*+\ell_{A(e_m)}^*\bigr).
\]
\end{lema}
\begin{proof}
La suma de probabilidades individuales se obtiene con el Lema~\ref{macro:casilla}. Al sustituir las
masas diagonales por sus objetivos se comete el error de \eqref{macro:predictor} dividido por $k$ y
multiplicado por el número de líneas. Cada intersección de dos líneas distintas tiene probabilidad a
lo sumo $\max_e p_e/k=O_\delta(k^{-2})$. Las dos primeras desigualdades de inclusión y exclusión
acotan el error de la unión por la suma de esas intersecciones.

Para una línea distinguida disponible $\lambda$, escribimos $Y_{e|\lambda}=\prod_{f\ni
e,\,f\ne\lambda}X_f$. Si $Y_{e|\lambda}=1$, su contribución disminuye cuando se alcanza alguna de las
otras tres líneas sin ocupar $\lambda$. Excluir la ocupación simultánea de $\lambda$ y de una de las
tres líneas cambia la probabilidad en $O_\delta(k^{-2})$. Si se ocupa $\lambda$, la regla de
congelación conserva el observable. Por tanto, con $\nabla_s$ denotando el incremento de un paso,
\[
 \E_{\rm B}[\nabla_s F_{\lambda,I}\mid\mathcal F_s]
 =-\frac1{kn}\sum_{e\in\lambda_I}Y_{e|\lambda}H_\lambda^*(e)
   +O_\delta\!\left(\frac{\zeta+k^{-1}}{k}F_{\lambda,I}\right).
\]
La derivada temporal de cada densidad de fibra en la tabla es $-H^*_\lambda/x$ veces esa densidad, por
\eqref{macro:derivadas}. El incremento de su objetivo sin congelación es la integral de esta derivada
dividida por $n$, con error de Taylor $O_\delta(n^{-2}|\lambda_I|)$, donde $|\lambda_I|$ mide el
intervalo de su parámetro. La probabilidad de congelación es $O(1/k)$ y el incremento ordinario del
objetivo es $O_\delta(|\lambda_I|/n)$. Congelar el objetivo en el paso en que se ocupa la línea añade
por tanto
$O_\delta(|\lambda_I|/(kn))$ a su deriva condicional. Una fibra ya congelada tiene ambas derivas
iguales a cero.

Para las masas de filas o columnas se tiene exactamente
$\E_{\rm B}[\nabla_s N_I\mid\mathcal F_s]=-N_I/k$ y el objetivo disminuye en $|I|/n$. Para el
rectángulo se aplica la fórmula de la unión a las cuatro líneas de filas y columnas y las $2|T|$
líneas diagonales. Son distintas para las cuádruplas ordinarias. Si $Y_{T}$ denota el producto de sus
indicadores,
\[
 \E_{\rm B}[\nabla_s P_{\Theta,T}\mid\mathcal F_s]
 =-\frac1{kn^4}\sum_{(i,i',j,j')\in\Theta}^{\rm ord}Y_T H_T^*
   +O_\delta\!\left(\frac{\zeta+k^{-1}}{k}P_{\Theta,T}\right).
\]
Cuando $Y_T=0$ su contribución sigue valiendo cero. La densidad del objetivo tiene derivada
exactamente igual a $-H_T^*/x$ veces ella misma. La segunda derivada temporal está uniformemente
acotada para $x\ge\delta$, de modo que el error de Taylor integrado es $O_\delta(|\Theta|/n^2)$.
\end{proof}

\begin{proposicion}[Cota de parada bajo la ley auxiliar]\label{macro:parada}
Sea $M=\lceil\delta n\rceil$. Para todo radio $\varrho_{\rm mac}>0$ suficientemente pequeño, se
puede fijar una partición $\mathcal I$ de paso $a_{\rm mac}$ y constantes
$c_{\delta,\varrho_{\rm mac}},C_{\delta,\varrho_{\rm mac}}>0$ de modo que, para $n$ suficientemente
grande,
\[
 P^{\rm B}\!\left(\max_{0\le s\le n-M}\Delta(s)\ge\varrho_{\rm mac}\right)
 \le C_{\delta,\varrho_{\rm mac}}\exp(-c_{\delta,\varrho_{\rm mac}}n).
\]
Se puede imponer el mismo control a una partición de observación fija pasando a un refinamiento suyo.
\end{proposicion}
\begin{proof}
Sea $\tau$ el primer índice con $\Delta(s)\ge\varrho_{\rm mac}$, limitado por $n-M$. Paramos en
$\tau$ las discrepancias y los incrementos centrados de cada observable menos su objetivo. El
incremento que llega a $\tau$ se incluye si la historia hasta el padre seguía dentro del radio
$\varrho_{\rm mac}$. Sea $\mathcal M^*$ la suma de los máximos temporales absolutos de esas
martingalas para intervalos de filas y columnas y para cajas y subconjuntos, más el máximo sobre las
líneas distinguidas de la suma de los máximos absolutos de sus intervalos.

En una caja o intervalo restamos en la deriva empírica y en la del objetivo un mismo valor de la tasa
$H^*$, tomado en un punto de esa región. El término de tasa constante es a lo sumo $C_\delta/n$ veces
la discrepancia absoluta de masa en esa región. La oscilación de $H^*$ es $O_\delta(a_{\rm mac})$ por
regularidad espacial y multiplica la suma de la masa empírica y la masa objetivo. En todas las
subfibras de una línea esas masas suman a lo sumo dos. Para cada uno de los dieciséis tipos, también
suman a lo sumo dos sobre todas las cajas. En los errores de Taylor las longitudes o volúmenes suman a
lo sumo uno por familia o tipo. Estos hechos dan, con una constante $C_\delta$ independiente del
número de intervalos,
\begin{equation}\label{macro:gronwall}
 \Delta(s\wedge\tau)\le \Delta(0)+\mathcal M^*
 +\frac{C_\delta}{n}\sum_{j<s\wedge\tau}\Delta(j)
 +C_\delta(a_{\rm mac}+C_{\mathcal I}/n),\qquad 0\le s\le n-M .
\end{equation}
Para las fibras primero se suman los errores absolutos sobre sus intervalos y luego se toma el máximo
sobre las líneas. La contribución de cada fibra cesa desde su congelación. Las estimaciones del
Lema~\ref{macro:deriva} sólo emplean en la deriva $\zeta$, que es lineal en $\Delta_0+\Delta_F$. El
exponente $1/4$ del Lema~\ref{macro:brecha} no afecta al coeficiente de \eqref{macro:gronwall}.

Hay $O_{\mathcal I}(n)$ martingalas escalares. Para una fibra que sigue disponible cada línea recién
ocupada afecta a lo sumo una casilla de la fibra, pues las líneas distintas se cortan a lo sumo una
vez. En el paso que ocupa la línea distinguida se congela toda su observación. Para un rectángulo,
requerir que una de las líneas ocupadas pase por una posición fija del rectángulo determina una
relación lineal no idéntica y permite $O(n^3)$ cuádruplas. Cada incremento de una masa normalizada,
incluido el del objetivo, tiene por tanto módulo a lo sumo $C_\delta/n$. La misma cota, con otra
constante, vale después de centrar y parar los incrementos.

Si $(\xi_j)$ son diferencias de martingala con $|\xi_j|\le d_j$, para $\theta>0$ la convexidad de la
exponencial entre $-d_j$ y $d_j$ y la media condicional cero dan
$\E[\exp(\theta\xi_j)\mid\mathcal F_{j-1}]\le\cosh(\theta d_j)\le\exp(\theta^2d_j^2/2)$. El producto
compensado es una supermartingala positiva. Pararla en el primer $h$ con $\sum_{j\le h}\xi_j\ge u$,
aplicar la misma cuenta a $-\xi_j$ y optimizar $\theta$ da
\begin{equation}\label{macro:max-mart}
 P\left(\max_{h\le m}\left|\sum_{j\le h}\xi_j\right|\ge u\right)
 \le 2\exp\!\left(-\frac{u^2}{2\sum_{j\le m}d_j^2}\right).
\end{equation}
Para la partición $\mathcal I$ fija existe $C'_{\mathcal I}$ tal que si todas las martingalas escalares
tienen máximo absoluto menor
que $b/C'_{\mathcal I}$ entonces $\mathcal M^*<b$. El factor $C'_{\mathcal I}$ cuenta los intervalos de
una sola fibra y las
observaciones de intervalos y rectángulos. El número de líneas distinguidas sólo interviene en la
unión de sucesos. De \eqref{macro:max-mart} resulta
\[
 P^{\rm B}(\mathcal M^*\ge b)\le C_{\mathcal I}(n+1)\exp(-c_{\delta,\mathcal I}b^2 n).
\]

Al inicio $x=U_f=1$. Las sumas de retícula por intervalos y cajas, la aproximación de los extremos
de las fibras y las $O(n^3)$ cuádruplas excepcionales dan $\Delta(0)\le C_{\mathcal I}/n$,
uniformemente en la línea distinguida. Elegimos primero $\varrho_{\rm mac}$ dentro de los umbrales
deterministas, después $a_{\rm mac}$ y $b$ con
$e^{C_\delta}(b+C_\delta a_{\rm mac})<\varrho_{\rm mac}/3$, y finalmente $n$ grande a partición
fija. Gronwall discreto en \eqref{macro:gronwall} da $\Delta(s\wedge\tau)<\varrho_{\rm mac}$ siempre
que $\mathcal M^*<b$, también para el paso que llega a $\tau$. Esto excluye la primera salida y
prueba la cota, absorbiendo $n+1$ en la exponencial. Al refinar una partición todas las masas y
objetivos de una región gruesa se obtienen sumando los de las regiones finas, también para fibras
congeladas. La desigualdad triangular transfiere la cota a la partición gruesa.
\end{proof}
